% Build from documentation/: latexmk -pdf -interaction=nonstopmode -halt-on-error sparsest_representation.tex
\documentclass[11pt,a4paper]{article}
\usepackage[T1]{fontenc}
\usepackage[utf8]{inputenc}
\usepackage{lmodern,microtype}
\usepackage[margin=24mm]{geometry}
\usepackage{amsmath,amssymb,amsthm,mathtools}
\usepackage{booktabs,tabularx,array,enumitem,fancyhdr,xcolor}
\usepackage[colorlinks,linkcolor=blue!45!black,citecolor=blue!45!black,urlcolor=blue!45!black,breaklinks]{hyperref}
\usepackage{bookmark}
\newtheorem{theorem}{Theorem}
\newtheorem{lemma}[theorem]{Lemma}
\newtheorem{corollary}[theorem]{Corollary}
\newtheorem{proposition}[theorem]{Proposition}
\theoremstyle{definition}
\newtheorem{definition}{Definition}
\newtheorem{example}{Example}
\theoremstyle{remark}
\newtheorem*{remark}{Remark}
\newcommand{\R}{\mathbb R}
\newcommand{\C}{\mathbb C}
\newcommand{\F}{\mathbb F}
\newcommand{\N}{\mathbb N}
\newcommand{\supp}{\operatorname{supp}}
\newcommand{\spanof}{\operatorname{span}}
\newcommand{\dist}{\operatorname{dist}}
\newcommand{\spark}{\operatorname{spark}}
\newcommand{\norm}[1]{\left\lVert#1\right\rVert}
\newcolumntype{Y}{>{\raggedright\arraybackslash}X}
\setlength{\parindent}{0pt}
\setlength{\parskip}{5pt}
\setlist{nosep,leftmargin=*}
\renewcommand{\arraystretch}{1.16}
\pagestyle{fancy}
\fancyhf{}
\fancyhead[L]{\small LRAC / Foundations of sparsest representation}
\fancyhead[R]{\small 29 October 2026}
\fancyfoot[C]{\thepage}
\setlength{\headheight}{14pt}
\hypersetup{pdftitle={Existence of Sparsest Representations: Theorems, Proofs, and Limits},pdfauthor={LRAC research project}}
\title{Existence of Sparsest Representations\\[5pt]\large Theorems, proofs, and the limits of what existence establishes}
\author{LRAC research project}
\date{29 October 2026}

\begin{document}
\maketitle
\begin{abstract}
We give a precise answer to whether every sequence has a sparsest possible
representation. First, any declared representation family with at least one
feasible finite code attains a minimum nonzero count, by well-ordering; compactness
is unnecessary for this conditional existence result. Second, every nonzero
sequence admits an exact one-coefficient representation in a sequence-dependent
algebraic basis. In finite dimensions this basis can be orthonormal and
constructed by Gram--Schmidt; the same conclusion holds for finite-energy infinite
sequences in a sequence-dependent Hilbert basis. These are existence statements,
not evidence of compression: the chosen basis can contain the sequence itself.
We then prove existence in fixed dictionaries, a conditional uniqueness result,
and limitations for shared stable decoders and finite-precision codes. No theorem
here asserts that the trained autoencoder attains a global optimum, or that one
universal, useful sparsifying domain has been identified.
\end{abstract}

\section{The question and its quantifiers}

A finite real recording is a vector $x\in\R^N$. More generally, a sequence is an
element of $\F^I$, where $\F=\R$ or $\C$ and $I$ is a finite index set or $\N$.
An infinite sequence need not belong to $\ell^2$; the space and reconstruction
criterion must therefore be stated explicitly.
We use $\N=\{1,2,\ldots\}$ and $\N_0=\{0,1,2,\ldots\}$; finite signal and
dictionary dimensions are positive integers.

The word \emph{sparsest} refers here to the \textbf{minimum number of nonzero
scalar coefficients}, not minimum bit length, smallest coefficient magnitudes,
or most interpretable features. We distinguish:
\begin{enumerate}
\item existence of a minimum in a declared feasible family;
\item existence when a new basis may be chosen for each individual sequence;
\item existence under a shared decoder, restricted adaptation and specified precision;
\item uniqueness, computation and certification of that minimum.
\end{enumerate}
These assertions are not interchangeable. In particular,
$\forall x\,\exists D_x$ does not imply $\exists D\,\forall x$ with the same sparse
budget. All results below are mathematical statements under their explicit
assumptions; they are not new empirical findings about speech.

\section{A general existence theorem}

\begin{definition}[Admissible finite-code representation]\label{def:admissible}
Fix a signal space $X$ and shared reconstruction conventions $h$. An admissible
representation is a pair $r=(\theta,z)$ in a declared set $\mathcal A_h$, where
$z\in\F^m$ for some finite $m\geq0$, and $\theta$ specifies any permitted
sample-dependent choice of transform, basis, decoder parameters or side
information. A fixed reconstruction rule $G_h$ maps $r$ to $X$.

Its coefficient cost is
\[
c(r)=\norm{z}_0=\#\{j:z_j\neq0\}\in\N_0.
\]
For a declared nonnegative distortion $d$ and tolerance $\varepsilon\geq0$, set
\[
\mathcal F_\varepsilon(x)
=\{r\in\mathcal A_h:d(x,G_h(r))\leq\varepsilon\}.
\]
Exact reconstruction can instead be required directly as $G_h(r)=x$.
Any restrictions on adaptation, coefficient range, precision, code length and
decoder access are part of $\mathcal A_h$ and $h$.
\end{definition}

\begin{lemma}[Attainment of an integer-valued minimum]\label{lem:integer}
Let $S$ be nonempty and $c:S\to\N_0$. There exists $s_*\in S$ such that
$c(s_*)\leq c(s)$ for all $s\in S$.
\end{lemma}
\begin{proof}
The image $c(S)$ is a nonempty subset of $\N_0$, so the well-ordering principle
gives a least element $k_*$. Since $k_*\in c(S)$, some $s_*\in S$ satisfies
$c(s_*)=k_*$. This $s_*$ has the required property.
\end{proof}

\begin{theorem}[Existence of a sparsest feasible finite representation]\label{thm:existence}
For every $x\in X$ and $\varepsilon\geq0$, if
$\mathcal F_\varepsilon(x)\neq\varnothing$, then there exists
$r_*=(\theta_*,z_*)\in\mathcal F_\varepsilon(x)$ satisfying
\[
k_*(x;\varepsilon,\mathcal A_h)
=\norm{z_*}_0
=\min_{(\theta,z)\in\mathcal F_\varepsilon(x)}\norm{z}_0.
\]
Consequently, any admissible feasible family for every signal in a class
$\mathcal X\subseteq X$ has a sparsest finite representation for every
$x\in\mathcal X$.
\end{theorem}
\begin{proof}
Apply Lemma~\ref{lem:integer} to
$S=\mathcal F_\varepsilon(x)$ and $c(\theta,z)=\norm{z}_0$.
\end{proof}

\begin{remark}[What is, and is not, assumed]
This result permits infinitely many transforms, noncompact parameter sets and
different finite code lengths. It needs no continuity, convexity or compactness:
the \emph{attained coefficient counts} are integers. It does require an actual
feasible finite code. It establishes neither uniqueness nor an algorithm for
finding the minimizer. The minimum distortion for a \emph{fixed} code budget is
a different optimization problem and can fail to be attained.
\end{remark}

\begin{example}[An infimum of error is not feasibility]\label{ex:infimum}
For scalar target $x=0$ and fixed decoder $G(z)=e^z$, $z\in\R$,
$\inf_z|x-G(z)|=0$, but $G(z)\neq0$ for every finite $z$. At tolerance zero there
is no feasible code. At every positive tolerance there is a feasible finite
code and Theorem~\ref{thm:existence} applies. Thus arbitrarily accurate
approximations do not, by themselves, prove exact feasibility.
\end{example}

For finite $x\in\F^N$, allowing the identity decoder with $z=x$ proves feasibility
with at most $N$ nonzeros. If that decoder is not admissible, feasibility must
be established separately. This is the first assumption to check for a learned
architecture with an information bottleneck.

\section{Unrestricted sequence-dependent bases}

\subsection{An algebraic statement for every sequence}

\begin{lemma}[Extension of a linearly independent set to a basis]\label{lem:extension}
Assume the axiom of choice in its equivalent Zorn-lemma form. If $V$ is a vector
space over $\F$ and $L\subseteq V$ is linearly independent, then $L$ is contained
in an algebraic (Hamel) basis of $V$.
\end{lemma}
\begin{proof}
Let $\mathcal P$ be the collection of linearly independent subsets of $V$
containing $L$, ordered by inclusion. It is nonempty since $L\in\mathcal P$.
For a nonempty chain in $\mathcal P$, its union contains $L$ and is linearly
independent: a finite linear relation in that union uses finitely many vectors,
all contained in one member of the chain. The empty chain has upper bound $L$.
Thus every chain has an upper bound, and Zorn's lemma gives a maximal element $B$.

If $\spanof(B)\neq V$, choose $v\notin\spanof(B)$. Then $B\cup\{v\}$ is linearly
independent, contradicting maximality. Hence $B$ spans $V$ and is a basis.
\end{proof}

\begin{theorem}[Exact global one-sparsity under unrestricted basis choice]\label{thm:one}
Let $V$ be a nonzero vector space over $\F$, including the space of length-$N$
finite sequences $\F^N$ or all infinite sequences $\F^\N$. Allow every algebraic basis
of $V$, with the basis known to the inverse. Then
\[
\min_{\substack{B\ \mathrm{basis\ of}\ V\\x=\sum_{b\in B}z_b b}}
\#\{b:z_b\neq0\}
=
\begin{cases}
0,&x=0,\\
1,&x\neq0,
\end{cases}
\]
where the sum is finite, as required for an algebraic basis.
\end{theorem}
\begin{proof}
For $x\neq0$, the singleton $\{x\}$ is linearly independent.
Lemma~\ref{lem:extension} supplies a basis $B_x$ containing $x$. Assign coefficient
$1$ to that basis vector and zero to every other vector; the reconstruction is
exact and the cost is one. No linear synthesis of an all-zero coefficient vector
can equal a nonzero $x$, so zero nonzeros are impossible. For $x=0$, all-zero
coefficients attain the lower bound zero.
\end{proof}

\begin{remark}[Meaning of the infinite-sequence assertion]
For all of $\F^\N$, this is an algebraic, generally nonconstructive statement.
A Hamel basis need not be countable, computable, continuous or stably encodable.
It is not a conventional Fourier/wavelet expansion or a finite-bit codec.
The theorem explicitly uses a basis that may depend on the target sequence.
\end{remark}

\subsection{Finite sequences: a constructive orthonormal result}

\begin{lemma}[Finite orthonormal completion]\label{lem:orthogonal}
For any unit vector $u\in\F^N$, there exists an orthonormal basis
$(u,q_2,\ldots,q_N)$ of $\F^N$.
\end{lemma}
\begin{proof}
Run Gram--Schmidt on the ordered spanning list
$(u,e_1,\ldots,e_N)$, discarding zero residuals. Its first output is $u$.
The process preserves the span of processed vectors and produces orthonormal
nonzero outputs. Since the original list spans $\F^N$, exactly $N$ outputs
are obtained, giving the required basis. This finite construction needs no
infinite-dimensional basis-extension argument.
\end{proof}

\begin{corollary}[Finite orthonormal sparsest-domain existence]\label{cor:orthonormal}
For every nonzero $x\in\F^N$, there is a unitary synthesis matrix $Q_x$ such that
\[
Q_x^*x=\norm{x}_2e_1,
\qquad
x=Q_x(\norm{x}_2e_1).
\]
The minimum nonzero count over all invertible linear coordinate systems, and
also over the smaller family of unitary coordinate systems, is exactly one.
For $x=0$, both minima are zero.
\end{corollary}
\begin{proof}
Apply Lemma~\ref{lem:orthogonal} to $u=x/\norm{x}_2$ and let $Q_x$ have the
resulting orthonormal vectors as columns. Inner products with these columns
give the displayed coefficient vector. The zero-code lower bound follows
from linearity, as in Theorem~\ref{thm:one}.
\end{proof}

\begin{corollary}[Unrestricted best sparsity at a norm-error tolerance]\label{cor:approx}
For $x\in\F^N$ and distortion $\norm{x-Qz}_2$, if all unitary $Q$ are allowed,
\[
k_*(x;\varepsilon)=
\begin{cases}
0,&\norm{x}_2\leq\varepsilon,\\
1,&\norm{x}_2>\varepsilon.
\end{cases}
\]
\end{corollary}
\begin{proof}
Zero coefficients reconstruct zero with error $\norm{x}_2$. If that error is
inadmissible, at least one coefficient is necessary, and
Corollary~\ref{cor:orthonormal} gives an exact one-coefficient reconstruction.
\end{proof}

Thus even orthogonality, norm preservation and condition number one do not stop
the basis from encoding the entire direction of $x$. These are not proposed
constraints for our model; they show that the issue persists under strong
linear regularity.

\subsection{Finite-energy infinite sequences}

\begin{theorem}[One-sparsity in a sequence-dependent Hilbert basis]\label{thm:hilbert}
For every nonzero $x\in\ell^2(\N;\F)$ there is a countable complete orthonormal
basis $(q_j)_{j\geq1}$ of $\ell^2$ such that
\[
q_1=x/\norm{x}_2,\qquad
x=\norm{x}_2q_1.
\]
Hence the exact minimum over these sequence-dependent Hilbert bases is one.
\end{theorem}
\begin{proof}
Apply Gram--Schmidt to $(x/\norm{x}_2,e_1,e_2,\ldots)$, skipping zero residuals.
The resulting orthonormal sequence has the same closed linear span as this
list. The list contains the canonical basis of $\ell^2$, whose finite linear
span is dense, so the resulting sequence is complete. Its first element is
$x/\norm{x}_2$, giving the claimed expansion. The zero-code lower bound is again
linearity.
\end{proof}

\begin{proposition}[Affine adaptation makes the minimum zero]\label{prop:affine}
If invertible affine coordinate systems may be chosen separately for each
$x\in\F^N$, every $x$ has an exact representation with zero nonzeros.
\end{proposition}
\begin{proof}
Choose $T_x(y)=y-x$, with inverse $T_x^{-1}(z)=z+x$. Then $T_x(x)=0$ and
$T_x^{-1}(0)=x$. Nonzero count is nonnegative, so zero is optimal.
\end{proof}

\textbf{Interpretation.} The linear theorem stores the sequence in its basis;
the affine theorem stores it in an offset. The inverse must know this
sample-dependent information. Neither proves economical reconstruction from
the counted coefficient vector alone.

\section{A nontrivial admissible family: a fixed finite dictionary}

Here the dictionary is fixed before the target is represented. These statements
do not optimize over every possible dictionary.

\begin{lemma}[Best approximation in a finite-dimensional subspace]\label{lem:projection}
Let $W\subseteq\F^N$ be a linear subspace. There is a unique $p_Wx\in W$ such that
\[
\norm{x-p_Wx}_2=\min_{y\in W}\norm{x-y}_2.
\]
\end{lemma}
\begin{proof}
Choose an orthonormal basis $(u_1,\ldots,u_r)$ of $W$ and put
$p_Wx=\sum_{j=1}^r\langle x,u_j\rangle u_j$, using the inner-product convention
linear in its first argument. The residual is orthogonal to $W$. For any $y\in W$,
\[
\norm{x-y}_2^2=\norm{x-p_Wx}_2^2+\norm{p_Wx-y}_2^2.
\]
This proves both minimality and uniqueness. For $W=\{0\}$, take $p_Wx=0$.
\end{proof}

\begin{theorem}[Existence and exact characterization in a fixed dictionary]\label{thm:dictionary}
Fix $D\in\F^{N\times M}$, $M\geq1$. For $S\subseteq\{1,\ldots,M\}$ write
$W_S=\spanof\{D_j:j\in S\}$, with $W_\varnothing=\{0\}$. Define
\[
\rho_k(x)=\min_{\substack{S\subseteq\{1,\ldots,M\}\\|S|\leq k}}
              \dist(x,W_S),\qquad 0\leq k\leq M .
\]
Each minimum is attained and equals
\[
\rho_k(x)=\min_{\norm{z}_0\leq k}\norm{x-Dz}_2.
\]
A feasible code at tolerance $\varepsilon$ exists exactly when
$\dist(x,\operatorname{range}D)\leq\varepsilon$. When feasible, a sparsest code
exists and its cost is
\[
k_*=\min\{k:\rho_k(x)\leq\varepsilon\}.
\]
In particular, if $D$ spans $\F^N$, every finite sequence has a sparsest exact
representation in $D$, with at most $N$ nonzeros.
\end{theorem}
\begin{proof}
For each support $S$, Lemma~\ref{lem:projection} gives an attained best waveform
$p_{W_S}x$. Since it belongs to the column span of $D_S$, it has coefficients
supported on $S$, even if these coefficients are nonunique. There are finitely
many supports with $|S|\leq k$, so a support attaining the smallest error exists.
Conversely, every $k$-sparse code synthesizes into one of these subspaces. This
proves the equality and attainment of $\rho_k$.

Taking all columns gives the feasibility criterion. Feasibility makes the
displayed set of integers nonempty, and its least element is attained by the
preceding construction. If $D$ spans $\F^N$, choose $N$ independent columns
forming a basis; they represent every $x$ exactly.
\end{proof}

No coefficient bound or compact parameter space is needed here: each support
produces a closed finite-dimensional waveform subspace, so its best waveform
approximation is attained. Exhaustive support search can be expensive; existence
is not a claim of efficient computation.

\begin{corollary}[Certified best basis in a finite orthonormal library]\label{cor:library}
Let $(Q_\ell)_{\ell=1}^J$ be a fixed finite library of unitary $N\times N$
matrices. For each $\ell$, let $a^{(\ell)}=Q_\ell^*x$, and order its magnitudes as
$|a^{(\ell)}|_{(1)}\geq\cdots\geq|a^{(\ell)}|_{(N)}$. Then
\[
\min_{\ell,\norm{z}_0\leq k}\norm{x-Q_\ell z}_2^2
=\min_{1\leq\ell\leq J}\ \sum_{j=k+1}^N|a^{(\ell)}|_{(j)}^2.
\]
The least $k$ for which the right side is at most $\varepsilon^2$ is the
globally sparsest feasible code in this declared library.
\end{corollary}
\begin{proof}
For fixed $\ell$, unitarity gives
$\norm{x-Q_\ell z}_2^2=\norm{a^{(\ell)}-z}_2^2$. On a chosen support the best
coefficients equal the corresponding entries of $a^{(\ell)}$. Keeping the
$k$ largest squared magnitudes minimizes the discarded sum. A finite comparison
over $\ell$ attains the best result.
\end{proof}

This supplies a rigorous small oracle for later experiments. It is a certificate
within the library, not across unknown transforms outside it. The library
identifier is part of the reconstruction information when the chosen basis is
not otherwise known to the inverse.

\begin{theorem}[Conditional uniqueness via spark]\label{thm:spark}
For a fixed finite dictionary $D$, let $\spark(D)$ be the smallest number of
linearly dependent columns, with convention $\spark(D)=M+1$ if all columns are
independent. If $Dz=x$ and
$2\norm{z}_0<\spark(D)$, then $z$ is the unique sparsest exact code in $D$.
\end{theorem}
\begin{proof}
Put $k=\norm{z}_0$. If another code $w\neq z$ satisfies $Dw=x$ and
$\norm{w}_0\leq k$, then $v=z-w\neq0$ satisfies $Dv=0$ and has at most $2k$
nonzeros. Its nonzero coordinates give a linear dependence among fewer than
$\spark(D)$ columns, a contradiction. Thus no distinct code of cost at most
$k$ exists.
\end{proof}

This standard result is included with its proof to separate \emph{existence}
from \emph{uniqueness}; see Donoho and Elad~\cite{donoho2003}. It is not a noisy
recovery theorem, a necessary uniqueness condition, or a certificate that $D$
is the best dictionary over all admissible choices.

\section{Why a shared stable decoder is a different claim}

\begin{lemma}[Lipschitz images of lower-dimensional domains have zero volume]\label{lem:null}
If $f:\R^s\to\R^N$ is Lipschitz and $0\leq s<N$, then
$f(\R^s)$ has zero $N$-dimensional Lebesgue measure.
\end{lemma}
\begin{proof}
For $s=0$, the image is a point and has zero volume. For $s\geq1$, restrict to
$[-R,R]^s$. Partition it into $n^s$ cubes of side $2R/n$. If $L>0$ is a
Lipschitz constant, the image of each cube has diameter at most
$2LR\sqrt{s}/n$, and fits in an $N$-dimensional box of side at most
$4LR\sqrt{s}/n$. The total volume of these covering boxes is bounded by
\[
n^s(4LR\sqrt{s}/n)^N=(4LR\sqrt{s})^N n^{s-N}\ \longrightarrow\ 0.
\]
The image therefore has zero outer measure. The case $L=0$ is constant.
Finally $\R^s=\bigcup_{R\in\N}[-R,R]^s$, and a countable union of null sets is null.
\end{proof}

\begin{theorem}[No universal exact low-support code under a shared Lipschitz decoder]\label{thm:uniform}
Fix $G:\R^M\to\R^N$ Lipschitz and an integer $0\leq k<N$. Then
\[
\mathcal Y_k=\{G(z):\norm{z}_0\leq k\}
\]
has zero $N$-dimensional measure and contains no nonempty open subset of $\R^N$.
Consequently, one such decoder cannot exactly reconstruct every waveform in an
open full-dimensional class using at most $k$ nonzeros.
\end{theorem}
\begin{proof}
For each support $S\subseteq\{1,\ldots,M\}$ with $|S|\leq k$, the coordinate
embedding $\iota_S:\R^{|S|}\to\R^M$ is an isometry. Thus $G\circ\iota_S$ is
Lipschitz and its image is null by Lemma~\ref{lem:null}. There are finitely many
supports. Their union is $\mathcal Y_k$, which is null. Every nonempty open
subset of $\R^N$ contains a box with positive volume, so none lies in $\mathcal Y_k$.
\end{proof}

A finite collection of such decoders also has a null union of $k$-sparse outputs.
If a single decoder receives $p$ additional continuous adaptation parameters
and is jointly Lipschitz, the same proof applies when $p+k<N$, using
$\R^p\times\R^{|S|}$ as the parameter domain. Adaptation cannot be ignored in
a dimension count.

This result does not contradict Theorem~\ref{thm:one}, whose basis may depend
freely on $x$. It also does not assert that speech contains an open subset of
waveform space. Nor does null volume by itself provide a positive
\emph{approximation} error bound: a null set may be dense. Approximate, stable
compression needs further code-range, precision and model-class assumptions.
Class-level stable approximation is treated more generally by
Cohen et al.~\cite{cohen2020}; their work is not a per-sequence shortest-code
certificate.

\section{Infinite sequences and finite-precision reconstruction}

\begin{example}[No finite exact code in a fixed infinite coordinate basis]\label{ex:infinite}
Let $x=(2^{-n})_{n\geq1}\in\ell^2$ and use its fixed canonical basis $(e_n)$.
Every coordinate of its unique exact expansion is nonzero, so no finite-support
exact representation exists in that basis. Theorem~\ref{thm:existence} does not
supply one because its exact finite-code feasible set is empty.

For support $S$, the best coefficient values are the original entries; the
captured energy is $\sum_{n\in S}4^{-n}$, maximized among $k$-element supports
by the first $k$ indices. Thus for any positive norm-error tolerance
$\varepsilon$, the best $k$-term canonical approximation has squared error
\[
\sum_{n=k+1}^\infty 4^{-n}=\frac{4^{-k}}{3}.
\]
Hence
\[
k_*=\max\!\left(0,\left\lceil\log_4\frac{1}{3\varepsilon^2}\right\rceil\right)
\]
is finite and attained. At $\varepsilon=0$ no finite $k$ suffices.
By Theorem~\ref{thm:hilbert}, the very same sequence is one-sparse in a suitably
chosen sequence-dependent Hilbert basis. The admissible family changes the answer.
\end{example}

\begin{theorem}[A finite-alphabet counting bound]\label{thm:counting}
Consider all length-$N$ sequences over an alphabet of size $q\geq2$. Suppose a
fixed reconstruction system accepts a choice from $J$ declared models and a code
$z$ of $M$ coordinates, each from an alphabet of $L\geq2$ coefficient symbols
including zero, with at most $k\leq M$ nonzeros. If it reconstructs all $q^N$
sequences exactly, it is necessary that
\[
q^N\ \leq\ J\sum_{j=0}^{k}\binom{M}{j}(L-1)^j .
\]
\end{theorem}
\begin{proof}
There are $\binom Mj$ choices of a support of size $j$ and $(L-1)^j$ assignments
of nonzero symbols on it. Multiplying the sum by $J$ bounds the total number of
possible input descriptions to the fixed decoder. A deterministic decoder
produces at most that many distinct outputs. To reconstruct $q^N$ distinct
sequences exactly, it must have at least $q^N$ distinct possible outputs.
\end{proof}

The support index and model choice are included in this bound. Increasing
coefficient precision or allowing uncounted sequence-specific models can evade a
coefficient-count comparison without reducing the actual information supplied.
This motivates explicit model/precision conventions, consistent with the
description-accounting perspective of Rissanen~\cite{rissanen1978}; it does not
replace minimum support with a different objective.

\begin{proposition}[No finite-bit exact description for every infinite binary sequence]\label{prop:binary}
Under any fixed decoder whose only input is a finite binary string, not every
infinite binary sequence has an exact description.
\end{proposition}
\begin{proof}
Finite binary strings form a countable set, so the decoder has at most countably
many outputs. Infinite binary sequences are uncountable: if $(s^{(j)})_{j\geq1}$
were an enumeration, the sequence $t_j=1-s^{(j)}_j$ would differ from its $j$th
listed sequence at coordinate $j$, contradicting completeness of the list.
Thus some infinite binary sequences are absent from the decoder's range.
\end{proof}

An algebraic one-coefficient representation from Theorem~\ref{thm:one} does not
contradict this proposition: its arbitrary basis and exact scalar need not have
finite binary descriptions. Likewise, mathematical finite real-valued vectors
and finite-precision stored recordings are distinct objects.

\section{What the existence proof establishes for this project}

The strongest unrestricted conclusion is already exact:
\begin{center}
\fbox{\parbox{0.91\linewidth}{
Every nonzero sequence admits an exact one-sparse representation in a
sequence-dependent algebraic basis. Every feasible finite-code problem has an
attained sparsest code within its declared admissible family. Neither statement
identifies a useful universal domain or an efficiently discoverable optimum.}}
\end{center}

For finite audio, our meaningful research target remains
\[
\min_{\theta,z}\norm{z}_0
\quad\text{subject to}\quad
d\bigl(x,G_h(\theta,z)\bigr)\leq\varepsilon,\qquad
(\theta,z)\in\mathcal A_h.
\]
The family may include unknown learned or sample-adapted domains. Its rules must
declare what the inverse knows, the allowed adaptation information, coefficient
precision/range, and any stability requirements. Existence does not force a
particular dictionary, orthogonality, semantic interpretation, or model architecture.

A feasible $k$-coefficient code gives an upper bound $U=k$ on the optimum.
A lower bound $L$ must rule out \emph{all} allowed codes and transform choices
with fewer coefficients. Matching bounds $L=U$ certify optimality within the
family. A unique sparsest code for one fixed dictionary does not prove that
dictionary is the best possible domain.

The frozen S-AF autoencoder is one admissible encoder--decoder construction.
Its successful reconstructions exhibit achievable codes. They do not prove
global sparsest-domain optimality. The next theoretical decision is the
admissible family and information-accounting rules; training can then be studied
as a search mechanism for that defined objective.

\begin{thebibliography}{9}
\bibitem{donoho2003}
D.~L. Donoho and M.~Elad.
\emph{Optimally sparse representation in general (nonorthogonal) dictionaries
via $\ell^1$ minimization}.
Proceedings of the National Academy of Sciences, 100(5):2197--2202, 2003.
\href{https://doi.org/10.1073/pnas.0437847100}{doi:10.1073/pnas.0437847100}.

\bibitem{cohen2020}
A.~Cohen, R.~DeVore, G.~Petrova and P.~Wojtaszczyk.
\emph{Optimal Stable Nonlinear Approximation}.
arXiv:2009.09907, 2020.
\url{https://arxiv.org/abs/2009.09907}.

\bibitem{rissanen1978}
J.~Rissanen.
\emph{Modeling by shortest data description}.
Automatica, 14(5):465--471, 1978.
\href{https://doi.org/10.1016/0005-1098(78)90005-5}{doi:10.1016/0005-1098(78)90005-5}.
\end{thebibliography}

\end{document}
